\documentclass[10pt,a4paper]{amsart}
\usepackage[margin=19mm]{geometry}
\usepackage{amsmath,amssymb,mathtools}
\usepackage[scr=rsfs,cal=euler]{mathalfa}
\usepackage{xfrac}
\usepackage[unicode,hidelinks,bookmarks=false]{hyperref}
\newcommand{\ie}{i.e.\ }
\newcommand{\cf}{cf.\ }
\newcommand{\resp}{respectively\ }
\newcommand{\Subsection}{\subsection}
\providecommand{\nobreakdash}{\nobreak\hbox{-}}
\setlength{\emergencystretch}{2em}
\multlinegap0pt
\usepackage[shortlabels]{enumitem}
\usepackage{amssymb}
\newtheorem{lemma}{Lemma}
\newcommand{\Xendd}{\operatorname{End}(X, \leq)}
\newcommand{\Z}{\mathbb Z}
\newcommand{\dom}{\operatorname{dom}}
\newcommand{\im}{\operatorname{im}}
\newcommand{\makeset}[2]{\left\lbrace #1 \; \colon\;
 \begin{tabular}{@{}l@{}}
   #2
  \end{tabular}
  \right\rbrace}
\title{Polish topologies on endomorphism monoids of linear orders}
\author{Serhii Bardyla, Luna Elliott}
\date{Source: arXiv:2605.26906v1 (CC BY 4.0)}
\begin{document}
\maketitle

\noindent\emph{Source and licence.} Polish topologies on endomorphism monoids of linear orders. Version arXiv:2605.26906v1 (CC BY 4.0), distributed by arXiv under the Creative Commons Attribution 4.0 International licence, http://creativecommons.org/licenses/by/4.0/, as recorded in the arXiv licence metadata for this exact version and shown on the abstract page https://arxiv.org/abs/2605.26906v1. Full licence notice: \texttt{licenses/arxiv-2605.26906-CC-BY-4.0.txt}.\par\smallskip

\noindent\textit{Editorial scope.} Counted unit: the article's lemma surjective (source0.tex lines 895--902). The article's complete original proof of it is delivered with it (source0.tex lines 903--975, one \texttt{proof} environment of 73 lines). No mathematics is written, repaired or re-derived: the statement and the proof are the article's own lines, in the article's own notation, and no retained line is substituted or re-flowed.  No further source-local context is needed: every cross-reference of every retained line resolves inside the two delivered spans.  Nothing was omitted from any retained span, so the package carries no omission record.  No retained line contains a citation, so the wrapper carries no bibliography and no bibliography entry.  No retained line contains a figure environment, an image-inclusion command or drawing code, so no image is delivered and the package declares no asset dependency.  Editorial formatting only, in addition: an AMS \texttt{amsart} preamble that restates the article's own declarations and macros used by the retained lines, namely the environments \texttt{lemma} on separate counters, together with the standard packages those lines need.  The retained environments print in this document as Lemma 1 in order of appearance, whereas the article prints the same items with the numbering of its own full document.  The complete native source of the article as distributed in the arXiv e-print for this exact version is bundled unchanged as \texttt{sources/201451/source0.tex}.
\begin{lemma}\label{strong surjective}
Let \(X\) be a  subset of \((\Z,\leq)\).
When giving \(\Xendd\) the pointwise topology, the following assertions hold: 
\begin{enumerate}
    \item for each surjective endomorphism $f$, the right shift $r_f$ is open;
    \item for each injective endomorphism $g$, the left shift $l_g$ is open.
\end{enumerate}
\end{lemma}

\begin{proof}
If $X$ is finite, then there is nothing to show as \(\Xendd\) is discrete, and thus each self-map of $\Xendd$ is open and closed. 

Assume that $X$ is infinite.
For all \(a,b\in X\) with \(a\leq b\) and \(f\in  \Xendd\), we define \(f_{a,b}\) to be the restriction of \(f\) to the set \(\makeset{x\in X}{\(a\leq x \leq b\)}\).
We then define
\[U^f_{a, b}:=\makeset{h\in \Xendd}{\(f_{a,b} \subseteq h\)}.\]
These sets form a basis for the topology on \(\Xendd\).

(1) Let \(f\in \Xendd\) be surjective, \(g\in \Xendd\) and $a,b\in X$ with $a\leq b$. 
In order to show that the right shift $r_f$ is open, it suffices to check that \(U_{a, b}^gf=U_{a,b}^{gf}\). The inclusion \(U_{a,b}^gf\subseteq U_{a,b}^{gf}\) is obvious. 
To show the converse inclusion fix any $w\in U_{a,b}^{gf}$. 
Note that for all \(k\in X\), the set \(I_k:=((k)w)f^{-1}\) is a nonempty convex subset of \(X\). Moreover the sets \(I_k\) are pairwise disjoint, and \(i_j<i_k\) whenever \(i_j\in I_j\), \(i_k\in I_k\) and \(j<k\). Since $(k)w=(k)gf$ providing \(a\leq k \leq b\), we have that \((k)g\in I_k\) for all \(a\leq k \leq b\).
For each \(k\in X\), the set \(\im(g_{a,b})\cap I_k\) is finite (possibly empty). So we can choose for each \(k\in X\) elements \(c^k_{\max}, c^k_{\min}\in I_k\) such that \(c^k_{\max}\) is not lower than any element of \(\im(g_{a,b})\cap I_k\), \((c^k_{\min}\) is no larger than any element of \(\im(g_{a,b})\cap I_k\), and \(c^k_{\max}\geq  c^k_{\min}\). Note that if $\im(g_{a,b})\cap I_k=\emptyset$, then  \(c^k_{\max}\geq  c^k_{\min}\) is the only condition  \(c^k_{\max}\) and \(c^k_{\min}\) are obliged to satisfy.
We define $h\in U_{a,b}^g$ as follows:
\[(k)h:=\begin{cases}
    c^k_{\min} & k< a\\
    c^k_{\max} & k> b\\
    (k)g_{a,b} & a\leq k\leq b.
\end{cases}\]
As \((k)h\in I_k\) for all \(k\in X\), it follows that \(hf=w\). 
In particular \(w=hf\in U^{g}_{a,b}f\) as required.

% the element of \(((k)w)f^{-1}\) with largest absolute value whenever \(|k|\geq n\). As \(w\in U^{gf}?n\), it follows that \(hf=w\). Moreocer 

% as follow
% \[(k)h=\begin{cases}
%     (k)g & |k|< n\\
%     ((m)w)f^{-1} &|k|\geq n
% \end{cases}\]

% $(k)h=(k)g$ for all $-n<k<n$, $(m)h=\max\big(((m)w)f^{-1}\big)$ for all $m\geq n$ and $(m)h=\min\big(((m)w)f^{-1}\big)$ for all $m\leq -n$. 

% Since the function $f$ is surjective, the map $h$ is well defined.
% It is straightforward to check that $hf=w$. We need to show that $h\in U_n^g$. Since $(k)h=(k)g$ for all \(-n<k<n\), it suffices to check that $h\in\Xendd$. Since $w,f\in\Xendd$, we get that  $$\max\big(((m)w)f^{-1}\big)\leq \max\big(((m+1)w)f^{-1}\big),$$ 
% $$\min\big(((m)w)f^{-1}\big)\geq \min\big(((m-1)w)f^{-1}\big),$$ 
% witnessing that $(m)h\leq (m+1)h$ for all $m$ such that $|m|\geq n$. 
% It remains to check that $(n)h\geq (n-1)h$ and $(-n)h\leq (-(n-1))h$. Taking into account that $(n-1)w=(n-1)hf$ and \((-(n-1))w=(-(n-1))hf\), we get $$(n)h=\max\big(((n)w)f^{-1}\big)\geq \max\big(((n-1)w)f^{-1}\big)=\max\big(((n-1)hf)f^{-1}\big)\geq (n-1)h,$$
% $$(-n)h=\min\big(((-n)w)f^{-1}\big)\leq \min\big(((-(n-1))w)f^{-1}\big)=\min\big(((-(n-1))hf)f^{-1}\big)\leq (n-1)h.$$
% Hence $h\in U_n^g$ and $U_n^gf=U_n^{gf}$, as required.

(2) Let $g, f\in\Xendd$ be such that \(g\) is an injection. We only need to show that the image of each neighborhood of \(f\) under \(l_g\) is a neighborhood of \(gf\).
Let \(a,b\in X\) be such that 
\begin{enumerate}
    \item \(a\leq b\),
    \item either \(a\in \im(g)\) or \(a=\min(X)\),
    \item either \(b\in \im(g)\) or \(b=\max(X)\).
\end{enumerate}
We show that \(gU^f_{a,b}\) is a neighborhood of \(gf\). This is sufficient as the family \[\{U_{a,b}^f: a,b \hbox{ satisfy conditions } (1)\hbox{--}(3)\}\] forms a neighbourhood base at \(f\).
Let \(a',b'\in X\) be such that whenever \(a\leq (k)g \leq b\), we have \(a'\leq k \leq b'\) (these exist as \(g\) is injective and there are only finitely many such \((k)g\)).
We show that \(U^{gf}_{a',b'} \subseteq gU^f_{a,b}\).
Let \(w\in U^{gf}_{a',b'}\) and consider the order preserving map \(g^{-1}w:\im(g) \to X\).
If \(h\in X^X\), then \(gh=w \iff h\supseteq g^{-1}w\). 
To conclude that \(w\in gU^f_{a,b}\), we need to show that \(g^{-1}w\cup f_{a,b}\) has an extension in \(\operatorname{End}(X)\).
It is easy to verify that every order preserving map between subsets of \(X\) has an extension to an element of \(\operatorname{End}(\Z,\leq)\) (in particular to an element of \(\operatorname{End}(X,\leq)\)).
Thus we only need to show that \(p:=g^{-1}w\cup f_{a,b}\) is an order preserving partial function.

We first show that \(p\) is a partial function. By the choice of \(a',b'\), if \(i\in \im(g)\) and \(a\leq i\leq b\) then \(a'\leq (i)g^{-1}\leq b'\). 
Thus \(((i)g^{-1})w=((i)g^{-1})gf=(i)f\) witnesses that \(g^{-1}w\) and \(f_{a,b}\) agree 
on $\dom(g^{-1}w)\cap \{x\in X: a\leq x\leq b\}=\im(g)\cap \{x\in X: a\leq x\leq b\}$.
Hence \(p\) is a partial function.

It remains only to show that \(p\) is order preserving. Let \(k,k^+\in \dom(p)\) be arbitrary such that \(k^+\) is the successor of \(k \) in the poset \(\dom(p)\). It suffices to show that \((k)p\leq (k^+)p\). There are 3 cases to consider:
\begin{enumerate}[\rm(a)]
    \item \(k<a\). As \(k^+\leq a\) and \(k^+\) is not minimal in \(X\), it follows from the choice of \(a\) (see condition (2)) that \(a\in \im(g)\).
    As \(k,k^+\in \dom(p)\), it follows that \(k,k^+\in \dom(g^{-1}w)=\im(g)\).
    \item \(a\leq k<b\). In this case \(k,k^+\in \dom(f_{a,b})\).
    \item \(b\leq k\). As \(k\geq b\) and \(k\) is not maximal in \(X\), it follows from the choice of \(b\) (see condition (3)) that \(b\in \im(g)\).
    As \(k,k^+\in \dom(p)\), it follows that \(k,k^+\in \dom(g^{-1}w)=\im(g)\).
\end{enumerate}

We showed above that in each case the points \(k,k^+\) belong to the domain of a common order preserving map contained in \(p\), which implies that \((k)p\leq (k^+)p\).
\end{proof}

\end{document}
